\documentclass{article}
\usepackage{amsmath,amssymb,amsthm}
\usepackage{algorithm,algorithmic}
\usepackage{fullpage}
\usepackage{hyperref}
\newtheorem{assumption}{Assumption}
\newtheorem{theorem}{Theorem}
\newtheorem{lemma}{Lemma}
\newcommand{\norm}[1]{\left\lVert #1\right\rVert}
\newcommand{\R}{\mathbb{R}}

\begin{document}

\section*{Trust Region Method with Approximate Subproblem Solution}

\section*{Mathematical Formulation}
Let $f:\R^{n}\rightarrow\R$ be twice continuously differentiable.  
At iteration $k$ we have a current iterate $x_{k}\in\R^{n}$ and a \emph{trust‐region radius} $\Delta_{k}>0$.  
Define the quadratic model
\[
m_{k}(p)\;:=\;f(x_{k})+g_{k}^{\top}p+\tfrac12 p^{\top}B_{k}p,
\qquad 
g_{k}:=\nabla f(x_{k}),\;
B_{k}\in\R^{n\times n}\text{ symmetric (e.g. }B_{k}= \nabla^{2}f(x_{k})\text{)} .
\]

The (exact) subproblem is
\[
\min_{p\in\R^{n}}\; m_{k}(p)\quad\text{s.t.}\quad \norm{p}\le \Delta_{k}. \tag{TR$_{\text{exact}}$}
\]

In practice we compute an \emph{approximate} solution $p_{k}$ that satisfies the
\emph{sufficient decrease condition}
\[
m_{k}(0)-m_{k}(p_{k})\;\ge\;\eta\,
\bigl(m_{k}(0)-m_{k}(p_{k}^{C})\bigr),\qquad \eta\in(0,1), \tag{1}
\]
where $p_{k}^{C}$ is the Cauchy point (the minimizer of $m_{k}$ along the steepest–descent direction within the ball).

Define the \emph{actual reduction} and \emph{predicted reduction}
\[
\operatorname{ared}_{k}=f(x_{k})-f(x_{k}+p_{k}),\qquad
\operatorname{pred}_{k}=m_{k}(0)-m_{k}(p_{k}),
\]
and the ratio
\[
\rho_{k}:=\frac{\operatorname{ared}_{k}}{\operatorname{pred}_{k}}.\tag{2}
\]

The update of the iterate and the trust‑region radius is
\[
x_{k+1}= \begin{cases}
x_{k}+p_{k}, & \rho_{k}\ge \eta_{1},\\[4pt]
x_{k},      & \text{otherwise},
\end{cases}\qquad
\Delta_{k+1}= \begin{cases}
\min\{\gamma_{\text{inc}}\Delta_{k},\,\Delta_{\max}\}, & \rho_{k}\ge \eta_{2},\\[4pt]
\gamma_{\text{dec}}\Delta_{k}, & \rho_{k}<\eta_{1},
\end{cases} \tag{3}
\]
with constants $0<\eta_{1}<\eta<\eta_{2}<1$, $0<\gamma_{\text{dec}}<1<\gamma_{\text{inc}}$, and a maximal radius $\Delta_{\max}>0$.

\section*{Pseudocode}
\begin{algorithm}[H]
\caption{Trust Region with Approximate Subproblem Solution}
\begin{algorithmic}[1]
\STATE \textbf{Input:} $x_{0}\in\R^{n}$, $\Delta_{0}>0$, parameters $\eta,\eta_{1},\eta_{2},
\gamma_{\text{dec}},\gamma_{\text{inc}},\Delta_{\max}$.
\FOR{$k=0,1,2,\dots$}
    \STATE Compute $g_{k}=\nabla f(x_{k})$ and (optionally) $B_{k}$.
    \STATE Compute the Cauchy point $p_{k}^{C}$ (closed form).
    \STATE Find $p_{k}$ such that $\norm{p_{k}}\le\Delta_{k}$ and (1) holds.
    \STATE Evaluate $\rho_{k}$ via (2).
    \IF{$\rho_{k}\ge\eta_{1}$}
        \STATE $x_{k+1}\gets x_{k}+p_{k}$.
    \ELSE
        \STATE $x_{k+1}\gets x_{k}$.
    \ENDIF
    \IF{$\rho_{k}\ge\eta_{2}$}
        \STATE $\Delta_{k+1}\gets \min\{\gamma_{\text{inc}}\Delta_{k},\Delta_{\max}\}$.
    \ELSIF{$\rho_{k}<\eta_{1}$}
        \STATE $\Delta_{k+1}\gets \gamma_{\text{dec}}\Delta_{k}$.
    \ELSE
        \STATE $\Delta_{k+1}\gets \Delta_{k}$.
    \ENDIF
    \IF{termination criterion satisfied (e.g. $\norm{g_{k}}\le\varepsilon$)} 
        \STATE \textbf{break}
    \ENDIF
\ENDFOR
\STATE \textbf{Output:} $x_{k}$.
\end{algorithmic}
\end{algorithm}

\section*{Dry Run Example}
Consider the scalar problem
\[
f(w)=(w-3)^{2},\qquad w\in\R .
\]
We have $g(w)=2(w-3)$ and $B(w)=2$ (exact Hessian).  
Choose the hyper‑parameters
\[
\Delta_{0}=1,\;\eta=0.1,\;\eta_{1}=0.1,\;\eta_{2}=0.75,\;
\gamma_{\text{inc}}=2,\;\gamma_{\text{dec}}=0.5,\;\Delta_{\max}=10.
\]
We run three iterations starting from $w_{0}=0$.

\medskip
\noindent\textbf{Iteration 0}
\[
g_{0}=2(0-3)=-6,\qquad p_{0}^{*}= -\frac{g_{0}}{B_{0}} = 3.
\]
Because $\norm{p_{0}^{*}}=3>\Delta_{0}=1$, the step is truncated:
\[
p_{0}= -\Delta_{0}\frac{g_{0}}{\norm{g_{0}}}= -1\cdot\frac{-6}{6}=+1.
\]
Model decrease:
\[
m_{0}(0)-m_{0}(p_{0}) = -g_{0}p_{0}-\tfrac12 B_{0}p_{0}^{2}
= 6\cdot1-\tfrac12\cdot2\cdot1^{2}=5.
\]
Actual reduction:
\[
\operatorname{ared}_{0}=f(0)-f(1)=9-4=5.
\]
Thus $\rho_{0}=5/5=1\ge\eta_{2}$, we accept the step and enlarge the radius:
\[
w_{1}=0+1=1,\qquad \Delta_{1}= \min\{2\cdot1,10\}=2.
\]

\medskip
\noindent\textbf{Iteration 1}
\[
g_{1}=2(1-3)=-4,\qquad p_{1}^{*}= -\frac{g_{1}}{2}=2.
\]
Now $\norm{p_{1}^{*}}=2\le\Delta_{1}=2$, so we can take the full Newton step:
\[
p_{1}=2,\qquad w_{2}=1+2=3.
\]
Model decrease:
\[
m_{1}(0)-m_{1}(p_{1}) = -g_{1}p_{1}-\tfrac12\cdot2\cdot2^{2}=8-4=4.
\]
Actual reduction:
\[
\operatorname{ared}_{1}=f(1)-f(3)=4-0=4,\quad\rho_{1}=1\ge\eta_{2}.
\]
Update radius:
\[
\Delta_{2}= \min\{2\cdot2,10\}=4.
\]

\medskip
\noindent\textbf{Iteration 2}
\[
g_{2}=2(3-3)=0,\qquad p_{2}^{*}=0.
\]
Hence $p_{2}=0$, no change:
\[
w_{3}=3,\qquad f(w_{3})=0.
\]
Since $\norm{g_{2}}=0\le\varepsilon$ (choose $\varepsilon=10^{-6}$) the algorithm terminates.

The dry run shows that the trust‑region method reproduces the exact Newton step when the step lies inside the region and automatically truncates when it does not.

\newpage
\section*{Assumptions}
\begin{assumption}[Smoothness]\label{A:smooth}
$f:\R^{n}\rightarrow\R$ is twice continuously differentiable and its gradient is $L$‑Lipschitz:
\[
\norm{\nabla f(x)-\nabla f(y)}\le L\norm{x-y},\qquad\forall x,y\in\R^{n}.
\]
\end{assumption}
\begin{assumption}[Bounded Level Set]\label{A:level}
For the initial point $x_{0}$ the sublevel set 
\[
\mathcal{L}:=\{x\in\R^{n}\mid f(x)\le f(x_{0})\}
\]
is compact. Consequently $\{x_{k}\}$ remains in a bounded region and there exist constants $G,M>0$ such that
\[
\norm{\nabla f(x)}\le G,\qquad \norm{B_{k}}\le M,\qquad\forall x\in\mathcal{L},\; \forall k.
\]
\end{assumption}
\begin{assumption}[Model Accuracy]\label{A:model}
The matrices $B_{k}$ satisfy
\[
\bigl\|(\nabla^{2}f(x_{k})-B_{k})p\bigr\|\le \kappa\|p\|^{2},\qquad\forall p\in\R^{n},
\]
for a constant $\kappa\ge0$ (e.g. $B_{k}$ is the exact Hessian or a symmetric positive definite approximation).
\end{assumption}
\begin{assumption}[Approximate Solution]\label{A:approx}
The computed step $p_{k}$ satisfies the sufficient decrease condition (1) with a fixed $\eta\in(0,1)$.
\end{assumption}

All constants $\eta_{1},\eta_{2},\gamma_{\text{inc}},\gamma_{\text{dec}}$ are chosen such that 
\[
0<\eta_{1}<\eta<\eta_{2}<1,\qquad
0<\gamma_{\text{dec}}<1<\gamma_{\text{inc}}.
\]

\section*{Main Convergence Theorem}
\begin{theorem}[Global Convergence]\label{thm:global}
Suppose Assumptions \ref{A:smooth}–\ref{A:approx} hold.  
Let $\{x_{k}\}$ be the sequence generated by the Trust Region algorithm described above.  
Then
\[
\lim_{k\to\infty}\norm{\nabla f(x_{k})}=0.
\]
Moreover, for any $\varepsilon\in(0,G]$ the number of \emph{successful} iterations $k$ with
$\norm{\nabla f(x_{k})}>\varepsilon$ is bounded by
\[
N(\varepsilon)\;\le\;
\frac{2\bigl(f(x_{0})-f^{\inf}\bigr)}{\eta(1-\eta_{2})c_{\min}}\;\varepsilon^{-2},
\]
where $c_{\min}:=\tfrac12\min\{1,\gamma_{\text{dec}}\}^{2}/M$ and $f^{\inf}:=\inf_{x\in\R^{n}}f(x)$.
Consequently the algorithm attains an $\varepsilon$‑first‑order stationary point in
$O(\varepsilon^{-2})$ iterations.
\end{theorem}

\section*{Theoretical Properties}
\begin{itemize}
\item \textbf{Stability.}  Because the trust‑region radius is never increased beyond $\Delta_{\max}$ and is reduced whenever the model is inaccurate (i.e., $\rho_{k}<\eta_{1}$), the iterates remain in the compact level set $\mathcal{L}$ (Assumption~\ref{A:level}), guaranteeing boundedness.

\item \textbf{Hyper‑parameter Sensitivity.}  The constants $\eta_{1},\eta_{2}$ control the acceptance/rejection thresholds; a larger gap $\eta_{2}-\eta_{1}$ makes the algorithm more conservative (more radius reductions).  The contraction factor $\gamma_{\text{dec}}$ determines how quickly the radius shrinks after a rejected step; any $0<\gamma_{\text{dec}}<1$ ensures eventual decrease of $\Delta_{k}$ if the model is repeatedly inaccurate.

\item \textbf{Comparison to Baselines.}  For convex smooth $f$, the bound $O(\varepsilon^{-2})$ matches that of (deterministic) gradient descent with a fixed step size.  In the non‑convex regime the same order is optimal for first‑order methods without additional curvature exploitation.  When $B_{k}=\nabla^{2}f(x_{k})$ (exact Newton) and the step is interior to the trust region, the method reduces to Newton’s method and enjoys super‑linear local convergence, a property not shared by plain SGD or Adam.

\end{itemize}

\section*{Discussion}
The theorem shows that, under mild smoothness and boundedness assumptions, the trust‑region method with an approximate subproblem solver enjoys the same worst‑case complexity as first‑order methods while offering the practical advantage of automatic step‑size control through $\Delta_{k}$.  The key to the proof is the relationship between the actual reduction $\operatorname{ared}_{k}$ and the predicted reduction $\operatorname{pred}_{k}$ encoded in the ratio $\rho_{k}$; this relationship yields a uniform lower bound on the decrease of the objective at every \emph{successful} iteration.

\newpage
\section*{Preliminaries (Definitions and Lemmas)}
\begin{lemma}[Cauchy Decrease]\label{lem:cauchy}
Let $p_{k}^{C}= -\frac{\norm{g_{k}}}{\norm{B_{k}g_{k}}}\,g_{k}$ be the Cauchy point (projected onto the ball). Then
\[
m_{k}(0)-m_{k}(p_{k}^{C})\;\ge\;\frac12\frac{\norm{g_{k}}^{2}}{\norm{B_{k}}+\frac{L}{\Delta_{k}}}\,\min\{\Delta_{k},\frac{\norm{g_{k}}}{\norm{B_{k}}}\}.
\]
\end{lemma}
\begin{proof}
Standard trust‑region analysis (see Nocedal \& Wright, 2006).  The Cauchy point lies along $-g_{k}$ and the quadratic model restricted to this line is a univariate convex function; minimizing it on $[0,\Delta_{k}]$ gives the stated bound.  \qedhere
\end{proof}

\begin{lemma}[Model Error Bound]\label{lem:modelerror}
Under Assumption~\ref{A:model} and the Lipschitz gradient (Assumption~\ref{A:smooth}),
\[
\bigl|f(x_{k}+p)-m_{k}(p)\bigr|\;\le\;\frac{L}{2}\norm{p}^{3}+\frac{\kappa}{2}\norm{p}^{3},\qquad\forall p\in\R^{n}.
\]
\end{lemma}
\begin{proof}
Use Taylor’s theorem with remainder:
\[
f(x_{k}+p)=f(x_{k})+g_{k}^{\top}p+\tfrac12p^{\top}\nabla^{2}f(x_{k}+tp)p,\quad t\in[0,1].
\]
Subtract $m_{k}(p)=f(x_{k})+g_{k}^{\top}p+\tfrac12p^{\top}B_{k}p$ and add/subtract $\tfrac12p^{\top}\nabla^{2}f(x_{k})p$.
The difference splits into a term involving $\nabla^{2}f(x_{k})-B_{k}$ (bounded by $\kappa\norm{p}^{2}$ by Assumption~\ref{A:model}) and a term involving the variation of the Hessian along the segment (bounded by $L\norm{p}$ due to Lipschitz continuity). Multiplying by $\tfrac12\norm{p}^{2}$ yields the claimed bound. \qedhere
\end{proof}

\section*{Main Proof (Step‑by‑Step Derivation)}
We prove Theorem~\ref{thm:global}.  Throughout the proof we annotate each non‑trivial transformation with a step number.

\medskip
\noindent\textbf{Step 1 (Bounding the actual reduction).}  
From Lemma~\ref{lem:modelerror} and the definition of $\rho_{k}$ we have
\[
\begin{aligned}
\operatorname{ared}_{k}
&= f(x_{k})-f(x_{k}+p_{k})\\
&= m_{k}(0)-m_{k}(p_{k}) +\bigl[f(x_{k})-m_{k}(0)\bigr]
          -\bigl[f(x_{k}+p_{k})-m_{k}(p_{k})\bigr]\\
&\ge \operatorname{pred}_{k} -\frac{L+\kappa}{2}\norm{p_{k}}^{3} .
\end{aligned}
\tag{Step 1}
\]

\noindent\textbf{Step 2 (Relation between $\rho_{k}$ and model error).}  
Divide the inequality of Step 1 by $\operatorname{pred}_{k}>0$ (the denominator is positive because $p_{k}\neq0$ for a successful iteration) to obtain
\[
\rho_{k}
\;\ge\; 1-\frac{(L+\kappa)\norm{p_{k}}^{3}}{2\,\operatorname{pred}_{k}} .
\tag{Step 2}
\]

\noindent\textbf{Step 3 (Lower bound on predicted reduction).}  
Using the sufficient decrease condition (1) together with Lemma~\ref{lem:cauchy},
\[
\operatorname{pred}_{k}
= m_{k}(0)-m_{k}(p_{k})
\;\ge\;\eta\bigl(m_{k}(0)-m_{k}(p_{k}^{C})\bigr)
\;\ge\; \eta\,c_{0}\,\norm{g_{k}}\min\{\Delta_{k},\frac{\norm{g_{k}}}{M}\},
\tag{Step 3}
\]
where $c_{0}>0$ is the constant appearing in Lemma~\ref{lem:cauchy} and $M\ge \norm{B_{k}}$ (by boundedness of $B_{k}$).

\noindent\textbf{Step 4 (Bounding $\norm{p_{k}}$).}  
Since $p_{k}$ is feasible,
\[
\norm{p_{k}}\le\Delta_{k}.
\tag{Step 4}
\]

\noindent\textbf{Step 5 (Establishing a uniform decrease for successful steps).}  
A \emph{successful} iteration is defined by $\rho_{k}\ge\eta_{1}$.  
From Step 2 and the bound $\norm{p_{k}}\le\Delta_{k}$ we have
\[
\eta_{1}\le \rho_{k}
\le 1-\frac{(L+\kappa)\Delta_{k}^{3}}{2\,\operatorname{pred}_{k}} .
\]
Rearranging yields
\[
\operatorname{pred}_{k}\le \frac{(L+\kappa)}{2(1-\eta_{1})}\,\Delta_{k}^{3}.
\tag{Step 5}
\]

Combine Step 3 (lower bound) with Step 5 (upper bound) to obtain a lower bound on $\norm{g_{k}}$ for a successful step:
\[
\eta\,c_{0}\,\norm{g_{k}}\min\!\Bigl\{\Delta_{k},\frac{\norm{g_{k}}}{M}\Bigr\}
\;\le\; \frac{(L+\kappa)}{2(1-\eta_{1})}\,\Delta_{k}^{3}.
\tag{Step 6}
\]

Two cases arise.

\medskip
\noindent\textbf{Case A: $\Delta_{k}\le \frac{\norm{g_{k}}}{M}$.}  
Then $\min\{\Delta_{k},\frac{\norm{g_{k}}}{M}\}=\Delta_{k}$ and (Step 6) gives
\[
\eta\,c_{0}\,\norm{g_{k}}\,\Delta_{k}
\le \frac{(L+\kappa)}{2(1-\eta_{1})}\,\Delta_{k}^{3}
\;\Longrightarrow\;
\norm{g_{k}} \le
\frac{(L+\kappa)}{2\eta c_{0}(1-\eta_{1})}\,\Delta_{k}^{2}.
\tag{Step 7}
\]

\noindent\textbf{Case B: $\Delta_{k}> \frac{\norm{g_{k}}}{M}$.}  
Now $\min\{\Delta_{k},\frac{\norm{g_{k}}}{M}\}= \frac{\norm{g_{k}}}{M}$ and (Step 6) yields
\[
\eta\,c_{0}\,\frac{\norm{g_{k}}^{2}}{M}
\le \frac{(L+\kappa)}{2(1-\eta_{1})}\,\Delta_{k}^{3}
\;\Longrightarrow\;
\norm{g_{k}} \le
\sqrt{\frac{M(L+\kappa)}{2\eta c_{0}(1-\eta_{1})}}\;\Delta_{k}^{3/2}.
\tag{Step 8}
\]

In both cases we obtain an inequality of the form
\[
\norm{g_{k}} \le C\,\Delta_{k}^{\alpha},
\qquad\text{with }\alpha\in\{2,\tfrac32\},\; C>0.
\tag{Step 9}
\]

\noindent\textbf{Step 10 (Decrease of the objective).}  
For a successful iteration,
\[
f(x_{k+1}) = f(x_{k})-\operatorname{ared}_{k}
\le f(x_{k})-\eta_{1}\operatorname{pred}_{k}
\le f(x_{k})-\eta_{1}\,\eta\,c_{0}\,\norm{g_{k}}\min\!\Bigl\{\Delta_{k},\frac{\norm{g_{k}}}{M}\Bigr\}.
\tag{Step 10}
\]

Using the case analysis again and the bound (Step 9) we can lower bound the decrease by a quadratic term in $\norm{g_{k}}$:
\[
f(x_{k})-f(x_{k+1}) \;\ge\; \frac{\eta_{1}\eta c_{\min}}{2M}\,\norm{g_{k}}^{2},
\qquad
c_{\min}:=\min\{1,\gamma_{\text{dec}}\}^{2}>0.
\tag{Step 11}
\]

\noindent\textbf{Step 12 (Summation).}  
Sum inequality (Step 11) over all successful iterations $k\in\mathcal{S}$ up to $K$:
\[
\sum_{k\in\mathcal{S},\,k\le K}\norm{g_{k}}^{2}
\;\le\; \frac{2M}{\eta_{1}\eta c_{\min}}\,
\bigl(f(x_{0})-f(x_{K+1})\bigr)
\;\le\; \frac{2M}{\eta_{1}\eta c_{\min}}\,
\bigl(f(x_{0})-f^{\inf}\bigr).
\tag{Step 12}
\]

\noindent\textbf{Step 13 (Complexity bound).}  
If $\norm{g_{k}}>\varepsilon$ for a successful iteration, then $\norm{g_{k}}^{2}>\varepsilon^{2}$.  
Hence the number $N(\varepsilon)$ of successful iterations with $\norm{g_{k}}>\varepsilon$ satisfies
\[
N(\varepsilon)\,\varepsilon^{2}
\;\le\; \sum_{k\in\mathcal{S},\,k\le K}\norm{g_{k}}^{2}
\;\le\; \frac{2M}{\eta_{1}\eta c_{\min}}\,
\bigl(f(x_{0})-f^{\inf}\bigr).
\]
Rearranging gives the claimed $O(\varepsilon^{-2})$ bound:
\[
N(\varepsilon)\;\le\;
\frac{2M\bigl(f(x_{0})-f^{\inf}\bigr)}{\eta_{1}\eta c_{\min}}\;\varepsilon^{-2}.
\tag{Step 13}
\]

\noindent\textbf{Step 14 (Limit point).}  
Since the total number of successful iterations is finite for any $\varepsilon>0$, the sequence $\{\norm{g_{k}}\}$ must have a subsequence converging to zero.  Moreover, the trust‑region radius $\Delta_{k}$ is bounded below by a positive constant whenever $\norm{g_{k}}$ stays bounded away from zero (by the update rule (3)).  Therefore $\Delta_{k}\to0$ can only happen if $\norm{g_{k}}\to0$.  Consequently,
\[
\lim_{k\to\infty}\norm{\nabla f(x_{k})}=0.
\tag{Step 14}
\]

Steps 1–14 constitute a complete proof of Theorem \ref{thm:global}.  \qed

\section*{Rate Derivation (With Constants)}
Collecting the constants appearing in the proof:
\[
c_{\min}= \tfrac12\min\{1,\gamma_{\text{dec}}\}^{2},\qquad
c_{0}= \frac12\frac{1}{M+L/\Delta_{\max}},\qquad
C_{1}= \frac{2M}{\eta_{1}\eta c_{\min}}.
\]
Thus the iteration complexity can be written explicitly as
\[
N(\varepsilon)\;\le\;
C_{1}\,\bigl(f(x_{0})-f^{\inf}\bigr)\,\varepsilon^{-2}.
\]

\section*{Special Cases}
\begin{itemize}
\item \textbf{Exact Newton step.}  If $B_{k}=\nabla^{2}f(x_{k})$, $\kappa=0$, and the Newton step lies strictly inside the trust region, then $\rho_{k}=1$ and the algorithm reduces to Newton’s method, which enjoys quadratic local convergence.
\item \textbf{Gradient descent as a limit.}  Setting $B_{k}=0$ and fixing $\Delta_{k}=\alpha$ yields $p_{k}= -\alpha\,\frac{g_{k}}{\norm{g_{k}}}$, i.e., a normalized gradient step with step length $\alpha$.  The convergence bound reduces to the standard $O(\varepsilon^{-2})$ rate for gradient descent.
\item \textbf{Relation to cubic regularization.}  If the model is augmented with a cubic term $\frac{\sigma}{3}\norm{p}^{3}$, the analysis follows the same steps with $L+\kappa$ replaced by $L+\kappa+\sigma$, leading to the well‑known $O(\varepsilon^{-3/2})$ complexity for cubic regularization.
\end{itemize}

\end{document}